\section{Algorithms for $\NAND_2$-avoiding $\confam$: Reduction to Implication}
\label{sec:redtoImpl}
%
In this section, we show that for any $\NAND_{2}$-avoiding constraint family $\confam$, we can reduce $\SAT(\confam)$ to \Implication. Specifically, we obtain the following theorem.
\begin{thm}\label{thm:redtoImpl}
Let $\confam$ be a $\NAND_2$-avoiding constraint family and let $T_\Impl(n,k)$ denote the optimal running time to solve \Implication. There is a constant $c_\confam$ and computable $f(k)$ such that we can solve $\SAT(\confam)$ in time 
$f(k)( T_\Impl(n,k) + n^{c_\confam})\log n$.
\end{thm}
%
Together with Theorem~\ref{thm:implalgo}, this gives an $f(k)n^{4\sqrt{k} + c_\confam}$-time algorithm for any $\NAND_2$-avoiding constraint family $\confam$.
%
To prove the above theorem, we prepare some notation and helpful facts. Let $\phi$ be an arbitrary formula.  For any assignment $a$, we call $a'$ a \emph{minimal satisfying extension} of $a$, if $a'$ satisfies $\phi$, $a\le a'$, and no other satisfying assignment $a''\notin \{a,a'\}$ fulfills $a\le a'' \le a'$. The following lemma shows that there are only $f(k)$ many minimal extensions of weight at most $k$, and these minimal extensions can be computed in time $f(k) n^{c}$ for some constant $c$ independent of $k$. Intuitively, this follows by using the bounded search tree technique over violated constraints, where the depth of the search tree is bounded by $k$ and each branching step has at most $r$ possibilities.
%
\begin{lem}[{\cite[Lemma 2.3]{BulatovM14}}]\label{lem:minextensions}
Let $\confam$ be a finite constraint family of bounded arity $r$. There is a constant $c_\confam'$ such that given any instance $\phi$ of $\SAT(\confam)$ and assignment $a$, there are at most $\Oh(r^k)$ minimal extensions of $a$ of weight $k$, and we can compute these extensions in time $\Oh(r^k n^{c'_\confam})$. 
\end{lem}
%
As an immediate useful consequence, we obtain that for our algorithmic results, we may assume without loss of generality that $\confam$ is $0$-valid, i.e., each $f\in \confam$ is satisfied by the all-zeroes assignment.
\begin{cor}[{see also \cite[Lemma 4.1]{Marx05}}]\label{cor:wlog0valid}
We can reduce any instance of $\SAT(\confam)$ with parameter~$k$ to $\Oh(r^k)$ many instances of $\SAT(\confam')$ with a parameter bounded by $k$, where $\confam'$ is the set of all $0$-valid $f'$ that are represented by $\confam$.
\end{cor}
\noindent 
By definition, if $\confam$ does not represent $\NAND_2$, then also $\confam'$ does not represent $\NAND_2$, and it remains to give an $f(k)(T_\Impl(n,k) + n^{c_\confam})\log n$-time algorithm for \emph{$0$-valid} $\NAND_2$-avoiding $\confam$.  
%
In the remainder of this section, we will use the graph formulation of the \Implication problem: We are given a directed graph $G=(V,E)$ and the task is to find a closed set $S$ (recall that $S$ is closed, if for all $(u,v)\in E$ we have that $u\in S$ implies $v\in S$) of size $k$. Recall that for any vertex set $S\subseteq V$, $D(S)$ denotes the set of descendants of any vertex $s \in S$ (including the vertices in $S$).
%
Our aim is the following: Given a formula $\phi$ of $\SAT(\confam)$, we give a randomized construction of an \Implication instance~$G$ such that 
\begin{enumerate}[label=(\roman*),nosep]
\item any closed set $S$ in $G$ corresponds to a satisfying assignment of $\phi$, and 
\item with large enough probability, $G$ contains a closed set of size $k$ if $\phi$ has a weight-$k$ solution.
\end{enumerate}
To this end, we let $V=\{x_1, \dots, x_n\}$ and recall that, for any set $S\subseteq V$, we let $a_S: [n]\to \{0,1\}$ denote a corresponding assignment with $a_S(i) = 1$ iff $x_i \in S$. From now on, we often synonymously speak of closed sets $S\subseteq V$ in $G$ and the corresponding assignment $a_S$ for $\phi$.
%
The rough outline is as follows: we start with the graph $G=(V,\emptyset)$, and try to repeatedly ``fix'' some closed set $S$ that violates $\phi$, by determining a (random) implication consistent with a minimal satisfying extension of $S$. The main insight is that if $\confam$ avoids $\NAND_2$, then it suffices to make sure that all sets $D(v)$ for $v\in V$ are satisfying and this will automatically ensure that every closed set is satisfying.
%
Let us formally describe the algorithm: 
\begin{enumerate}
\item Given $\phi$, initialize $G= (V,E)$ with $V=\{x_1,\dots, x_n\}$ and $E=\emptyset$.
\item While there exists some $v\in V$ such that $a_{D(v)}$ violates $\phi$, do the following:
\begin{enumerate}
\item Compute the set $A_v$ of minimal satisfying extensions of $a_{D(v)}$ of weight at most $k$.
\item Let $X$ consist of all $x_i \in V\setminus D(v)$ such that there is some $a\in A_v$ with $a(i) = 1$.
\item If $X = \emptyset$, delete all ascendants of $v$ (including $v$) from $G$. Otherwise, pick $x$ uniformly at random from $X$ and add the edge $(v,x)$ to $E$.
\end{enumerate}
\end{enumerate}
%
The important properties of the algorithm are captured in the following lemma.
\begin{lem}\label{lem:randprop}
Let $\confam$ be a finite $0$-valid constraint family. There is a constant $c_\confam$ and a function $g(k)$ such that the following properties hold.
\begin{enumerate}[label=(P\arabic*)]
\item\label{item:iters} During the process, each vertex $v$ is considered at most $k$ times in the while loop. Thus, the algorithm can be implemented to run in time $\Oh(g(k)n^{c_\confam})$.
\item\label{item:satisclosedwhp} If $\phi$ has a satisfying assignment of weight $k$, then with probability at least $g(k)^{-1}$, there is a closed set $S$ in $G$ of size $k$.
\item\label{item:closedissat} If $\confam$ avoids $\NAND_2$, any closed set $S\subseteq V$ in the constructed graph yields a satisfying assignment $a_S$ for $\phi$. 
\end{enumerate}
\end{lem}
\begin{proof}
For \ref{item:iters}, note that whenever $v\in V$ is considered in the while loop, it is either deleted, or an edge $(v,x)$ with $x\notin D(v)$ is added to the graph. Thus, when $v$ is considered for the $k$-th time, we have $|D(v)|\ge k$, and thus there can be no satisfying extension of $a_{D(v)}$ of weight at most $k$. Consequently, we must have $A_v=\emptyset$, and thus $X=\emptyset$, which forces $v$ to be deleted. Thus, we have at most $kn$ iterations of the while loop, where each iteration can be implemented in time $\Oh(r^k n^{c'_\confam})$ by Lemma~\ref{lem:minextensions}.
%
For~\ref{item:satisclosedwhp}, assume that there is a set $S$ of size $k$ such that $a_S$ satisfies $\phi$. We show that with large enough probability, we will maintain as invariant that $D(v)\subseteq S$ for every $v\in S$, and thus $S$ will be a closed set in $G$. To this end, we first observe that for $D(v) \subseteq S$ to hold for all $v\in S$, it suffices that the following property holds:
\begin{equation}\label{eq:processgood}
\text{In each iteration that considers a vertex $v\in S$, the selected vertex $x$ is in $S$.}
\end{equation}
Indeed, if this is the case, then no $v\in S$ is ever deleted. Furthermore, we have that $D(v)\subseteq S$ for all $v\in S$, and thus $S$ is a closed set in $G$. It remains to give a lower bound on the probability that~\eqref{eq:processgood} holds throughout the process.
%
To this end, consider the event that some $v\in V$ is considered in the while loop, conditioned that~\eqref{eq:processgood} has not been violated in a previous iteration. Under this event, $D(v) \subseteq S$, and thus there is a minimal satisfying extension $D(v) \subsetneq S' \subseteq S$ such that $a_{S'}$ satisfies $\phi$ and thus $a_{S'}\in A_v$. Let $s\in S'\setminus D(v)$ be arbitrary, then $s\in X$ by construction (note that $s$ has not been deleted). By Lemma~\ref{lem:minextensions}, we have that $|A_v| \le \Oh(r^k)$. Since each $a\in A_v$ has weight at most $k$, this yields $|X|\le k|A_v|\le \Oh(k r^k)$. Thus, the probability that the random choice is $x=s$ is at least $1/|X|\ge \Omega(1/(kr^k))$. Finally, we observe that by~\ref{item:iters}, for each $v\in S$, there are at most $k$ iterations considering $v$, where each iteration has a probability of at least $\Omega(1/(kr^k))$ of not violating~\eqref{eq:processgood}. Thus, we obtain that~\eqref{eq:processgood} holds with probability at least $\Omega(1/(kr^k)^{k|S|}) = \Omega(1/(kr^k)^{k^2})$, and the claim follows by setting $g(k) \coloneqq (kr^k)^{-k^2}$.
%
Finally, for \ref{item:closedissat}, note that at the end of the process, the property holds that
\begin{equation}\label{eq:Nvclaim}
\text{For all (remaining) } v\in V, a_{D(v)} \text{ satisfies } \phi.
\end{equation}
We will leverage this fact to show that $a_S$ satisfies $\phi$ for \emph{all} closed sets $S=D(v_1)\cup ... \cup D(v_\ell)$ for $v_1,\dots, v_\ell\in V$. We first transform the graph $G$ to a DAG by contracting all strongly connected components $C=\{v_1, \dots, v_c\}$ to a single vertex $v_C$ representing the set $C$. Note that the closed sets in the DAG remain in a one-to-one correspondence to the closed sets of the original graph (and the corresponding assignments to $\phi$), thus this transformation is without loss of generality. Thus, we may assume that $G$ has a topological ordering $v_1,\dots, v_{n'}$ of its vertices ($n'\le n$). We will prove by induction on $i=n', ..., 1$ that for all closed sets $S\subseteq \{v_i,...,v_{n'}\}$, $a_S$ satisfies $\phi$.
%
For the base case $i=n'$, we only need to verify that (i) the all-0 assignment satisfies $\phi$, which holds by 0-validity of $\confam$, and (ii) that $a_{v_{n'}}$ satisfies $\phi$, which holds by~\eqref{eq:Nvclaim} (as $D(v_{n'}) = \{v_{n'}\}$). Thus, for $i<n'$, let us assume that the claim holds for $i+1$. Consider any closed set $U\subseteq \{v_i, \dots, v_{n'}\}$. If $U$ does not contain $v_i$, the claim follows by inductive assumption, thus let us assume that $v_i \in U$ and thus $U\supseteq D(v_i)$, as $U$ is closed. If $U = D(v_i)$, $a_U$ satisfies $\phi$ by~\eqref{eq:Nvclaim}. Thus, it remains to consider $U \supsetneq D(v_i)$, for which we assume for contradiction that $a_U$ violates~$\phi$. Let $W := U\setminus D(v_i)$, and note that $D(W)\subseteq U$ is a closed set in $\{v_{i+1},\dots, v_{n'}\}$. Thus, by inductive assumption, $a_{D(W)}$ satisfies $\phi$. Furthermore, observe that $Z:=D(v_i)\cap D(W)$ is a closed set in $\{v_{i+1}, \dots, v_{n'}\}$ (since the intersection of any two closed sets yields a closed set). Thus, $a_Z$ satisfies $\phi$ by inductive assumption. It remains to show that the fact that $a_{D(v_i)}$, $a_{D(W)}$ and $a_Z=a_{D(v_i)\cap D(W)}$ all satisfy $\phi$, while $a_U=a_{D(v_i)\cup D(W)}$ violates $\phi$, gives a contradiction to $\confam$ avoiding $\NAND_2$.
%
To this end, let $C$ be a constraint violated by $a_U$ and note that $C=f(x_{i_1},\dots,x_{i_r})$ for some $f\in \confam$ and $i_1,\dots, i_r\in [n]$. Note that we can view $f$ as $f: \{0,1\}^{V_c} \to \{0,1\}$ for some appropriate variable set $V_C$. We show how to obtain $\NAND_2$ as a restriction of $f$ by partitioning $V_C$ into $X'\coloneqq (D(v_i) \setminus Z) \cap V_C, Y' \coloneqq (D(W)\setminus Z)\cap V_C, Z_1 \coloneqq Z\cap V_C, Z_0 \coloneqq V_C \setminus (X'\cup Y' \cup Z_1)$ and observing that 
\begin{align*}
\begin{matrix}
f( \overbrace{\entry{0}}^{X'},& \overbrace{\entry{0}}^{Y'}, & \overbrace{\entry{0}}^{Z_0}, & \overbrace{\entry{1}}^{Z_1}) & = & 1, & \text{[since $a_Z$ satisfies $C$]} \\
f( \entry{1},& \entry{0}, & \entry{0}, & \entry{1}) & = & 1, & \text{[since $a_{D(v_i)}$ satisfies $C$]} \\
f( \entry{0},& \entry{1}, & \entry{0}, & \entry{1}) & = & 1, & \text{[since $a_{D(W)}$ satisfies $C$]} \\
  \hline
f( \entry{1},& \entry{1}, & \entry{0}, & \entry{1}) & = & 0. & \text{[since $a_{D(W)\cup D(v_i)}$ violates $C$]} \\
\end{matrix}
\end{align*}
%
%
%
\end{proof}
%
It remains to give the proof of Theorem~\ref{thm:redtoImpl}.
\begin{proof}[Proof of Theorem~\ref{thm:redtoImpl}]
By Corollary~\ref{cor:wlog0valid}, we may assume without loss of generality that $\confam$ is 0-valid.
We repeat the following process $g(k)$ many times: We use the above algorithm to generate an \Implication instance $G$, and return YES if $G$ contains a closed set of size $k$, which we determine using an optimal \Implication algorithm. If none of the $g(k)$ iterations were successful, we return NO. Note that this approach can be implemented in time $g(k)\Oh(g(k)n^{c_\confam} + T_\Impl(n,k))$ by~\ref{item:iters}, and correctly decides the instance with probability at least $1-(1-1/g(k))^{g(k)} \ge 1-1/e$ by~\ref{item:satisclosedwhp} and~\ref{item:closedissat}.
%
The algorithm described above can be derandomized using the standard technique of Color Coding \cite{AlonYZ95}. In each iteration when vertex $v$ is considered, a random vertex $x$ is selected from a set $X$ of at most $K=\Oh(kr^k)$ vertices. As each vertex is considered at most $k$ times, we can represent the random choices by a function $r:V\to [K]^k$,  with the meaning that $r(v)$ is the vector of choices made when considering vertex $v$.  As discussed in the proof of Lemma~\ref{lem:randprop}, when considering vertices $v\in S$, these random choices need to be consistent with $S$ to ensure that $S$ is a closed set in the resulting graph. That is, for each $v\in S$ there is a vector $c(v)\in [K]^k$ such that if the random choice satisfies $r(v)=c(v)$ for every $v\in S$, then $S$ is a closed set.
%
We say that a family $\mathcal{H}$ of functions $h:[n]\to [k]$ is a {\em $(n,k)$-perfect family of hash functions} if for every $S\subseteq V$ of size $k$, there is an $h\in \mathcal{H}$ that is injective on $S$, i.e., assigns different values to different elements of $S$. It is known that a $(n,k)$-perfect family of size $2^{O(k)}\log n$ can be computed in time $2^{O(k)}n\log n$ \cite{AlonYZ95}. The derandomized algorithm would first compute such a family $\mathcal{H}$ over $V$ and would iteratively go through every $h\in \mathcal{H}$ and function $q:[k]\to [K]^k$. For a given choice of $h$ and $q$, we define the function $r(v)=q(h(v))$ and run the randomized algorithm using this function $r$ instead of the random choices. It is easy to see that the definition of $(n,k)$-perfect hash functions implies that there is at least one choice of $h$ and $q$ where $r(v)$ is exactly the prescribed value $c(v)$ for every $v\in S$ and therefore the randomized algorithm correctly finds the solution $S$. As we are considering at most $|\mathcal{H}|=2^{O(k)}\log n$ functions $h$ and $K^{k^2}$ different functions $q$, there is a function $f(k)$ such that the total running time is at most $f(k)\log n$ times a single run of the randomized algorithm.
\end{proof}
%
%
%
%
\section{Algorithms for $\NAND$-representing $\confam$: Reduction to Clique}
\label{sec:redToClique}
%
In this section, we develop algorithm for constraint families that might represent $\NAND_2$, but avoid $\NAND_d$ for some $d\ge 3$. To this end, we give a reduction to $(d-1)$-uniform HyperClique for $\NAND_d$-avoiding families, giving in particular a $f(k)n^{(\omega/3)k+c_\confam}$-time algorithm for $\NAND_3$-avoiding families. 
%
We first start with a natural reduction of $\SAT(\confam)$ for any $\confam$ with arity bounded by $r$ to $r$-uniform HyperClique, based on color-coding. To this end, let $\THC{d}(n,k)$ denote the optimal running time of finding a $k$-clique in a $d$-uniform hypergraph.
\begin{prop}\label{prop:aritybaseline}
Let $\confam$ be a constraint family of arity at most $r$. Then $\SAT(\confam)$ can be solved in time $f(k)(n^{2r} + \THC{r}(n,k))\log n$.
\end{prop}
\begin{proof}
Let $\phi$ be an arbitrary $\SAT(\confam)$ formula. Observe that any constraint $C$ of $\phi$ depends only on a set $\vars(C)\subseteq \{x_1,\dots,x_n\}$ of at most $r$ variables. For an assignment $a$, we let $C(a)\in \{0,1\}$ denote whether $C$ is satisfied by $a$.
%
We first show how to determine, given a partition of $x_1, \dots, x_n$ into $k$ sets $X_1,\dots, X_k$, whether there is a solution that sets precisely one variable in each $X_i$ to true. To this end, we construct a hypergraph $G$ with vertex set $X_1 \cup X_2 \cup \cdots \cup X_k$ and the following set of hyperedges: we include each possible hyperedge $e=\{x_{j_1}, \dots, x_{j_r}\}$ with $x_{j_1}\in X_{j_1}, \dots, x_{j_r} \in X_{j_r}$ and distinct $j_1, \dots, j_r\in [k]$ \emph{unless} there exists a clause $C$ with $\vars(C) \subseteq X_{j_1} \cup \cdots \cup X_{j_r}$ which is violated by the assignment that sets precisely the variables $e=\{x_{j_1}, \dots, x_{j_r}\}$ to $1$, i.e., $C(a_{e}) = 0$.
%
We claim that $H := \{x_{i_1}, \dots, x_{i_{k}}\}$ with  $x_{i_1} \in X_1, \dots, x_{i_k} \in X_k$ yields a $k$-clique in $G$ if and only if the assignment $a_{H}$ satisfies $\phi$. Indeed, assume that there is a clause $C$ violated by $a_{H}$. Note that as $C$ has arity at most $r$, we have $\vars(C) \subseteq X_{j_1} \cup \cdots \cup X_{j_r}$ for some distinct $i_1, \dots, i_r \in [k]$ (if $C$ involves variables of less than $r$ sets, we may use arbitrary additional sets). Thus, $e:=\{x_{j_1}, \dots, x_{j_r}\}$ cannot be an edge in $G$, since $a_H$ violates $C$, $a_e$ and $a_H$ agree on $\vars(C)$, and thus also $a_{e}$ violates~$C$. Conversely, if there is some $e:= \{x_{i_1}, \dots, x_{i_r}\}$ with distinct $i_1, \dots, i_r \in [k]$ such that $e$ is not an edge in $G$, then there exists some clause $C$ with $\vars(C) \subseteq X_{i_1} \cup \dots \cup X_{i_r}$ which is violated by $a_e$. Since $a_H$ and $a_e$ agree on $\vars(C)$, we conclude that also $a_H$ violates $C$ and thus $\phi$.
%
To create the desired $k$-partition of variables, we use a (deterministic) color-coding scheme: Let $\mathcal{H}$ be a \emph{$(n,k)$-perfect family of hash functions} $h : [n]\to [k]$ -- recall that this means that for any $S=\{s_1,\dots, s_k\}\subseteq [n]$, there exists some $h\in \mathcal{H}$ such that $\{h(s_1),\dots, h(s_k)\} = \{1,\dots, k\}$. Known efficient constructions~\cite{SchmidtS90, AlonYZ95} produce such assignments with $\ell = 2^{\Oh(k)}\log(n)$ in time $2^{\Oh(k)}n \log n$. Given this family, we create for each $h \in \mathcal{H}$ the $k$-partition $X^{(h)}_1, \dots, X^{(h)}_k$ with $X^{(h)}_j = \{ x_{s} \mid h(s) = j\}$ and solve the corresponding $r$-uniform HyperClique instance in time $\THC{r}(n,k)$. If any of these instances returns a solution, then indeed $\phi$ has a satisfiable assignment of weight $k$. Conversely, if $a_S$ is a weight-$k$ satisfying assignment for $\phi$, then by construction, there exists a hash function $h \in \mathcal{H}$ such that $|S\cap X^{(h)}_j|=1$ for $j=1,\dots, k$, and thus the corresponding $r$-uniform HyperClique instance indeed contains a solution. For each of the $2^{\Oh(k)}\log(n)$ hash functions, the time to construct and solve the $d$-uniform HyperClique instance is bounded by $\Oh(n^{2r} + \THC{r}(n,k))$, concluding the claim. 
\end{proof}
%
The main result in this section is the following reduction from $\NAND_{d+1}$-avoiding constraint families to $d$-uniform HyperClique.
%
\begin{thm}\label{thm:NANDavoidingalgo}
Let $d\ge 2$ and $\confam$ be an $\NAND_{d+1}$-avoiding constraint family. If there are constants $\gamma \ge d/(d+1)$ and $c$, and a computable $g(k)$ such that $d$-uniform HyperClique can be solved in time $g(k) n^{\gamma k + c}$, then there is a constant $c'$ and computable $g'(k)$ such that $\SAT(\confam)$ can be solved in time $g'(k) n^{\gamma k + c'}$.
\end{thm}
In particular, since we can find $k$-cliques in graphs in time $\Oh(n^{\frac{\omega}{3}k + 1})$, we obtain an $g(k) n^{\frac{\omega}{3}k + c'}$-time algorithm for solving $\SAT(\confam)$ for all $\NAND_3$-avoiding constraint families. Similarly, if for $d\ge 3$ the $d$-uniform HyperClique conjecture is refuted by exhibiting a $g(k)n^{(1-\eps)k+c}$-time algorithm for some constants $0< \eps < 1/(d+1)$ and $c$, we would obtain a $g'(k) n^{(1-\eps)k+c'}$-time algorithm for $\SAT(\confam)$ for $\NAND_{d+1}$-avoiding families $\confam$.
%
In the remainder of the section, we give the proof of Theorem~\ref{thm:NANDavoidingalgo}. 
The main task of the algorithm is to detect \emph{robust} assignments, defined as follows.
%
\begin{defn}
Let $a: [n] \to \{0,1\}$ be a weight-$k$ assignment that satisfies $\phi$. We say that $a$ is \emph{$d$-robust} if there is no assignment $a' \le a$ of weight at most $d$ that violates $\phi$.
\end{defn}
%
The first step of the algorithm is the easier task of detecting satisfying assignments that are \emph{not} $d$-robust (if there exists any):
Intuitively, an assignment that is not $d$-robust offers an advantage to find it: Assume we correctly guess an assignment $a'\le a$ of weight $w \le d$ such that some clause $C$ is violated by $a'$, then to extend $a'$ to the satisfying assignment $a$, we know that at least one additional variable in $C$ must be set to true. By bruteforcing over the at most $r-w \le r$ many possibilities, we gain an advantage. Specifically, by enumerating $O(n^{w}r) = \Oh(n^{w})$ many possibilities, we can fix $w+1$ true variables in our solution.
%
Let $T(n,k)$ denote the time our algorithms takes to solve an arbitrary $\SAT(\confam)$ instance for a $\NAND_{d+1}$-avoiding family $\confam$. 
In a preprocessing step, we first enumerate all assignments $a'$ of weight at most $d$. If there exists a clause $C_{a'}$ that is violated by $a'$, then we enumerate all variables $x \in \vars(C_{a'})\setminus \ones(a')$ (recall that $\vars(C)$ is the set of variables involved in $C$ and $\ones(a)$ denotes the set of variables set to $1$ under $a$). We recursively determine satisfiability of the formula $\phi_{a',x}$ obtained by restricting all variables in $\ones(a') \cup \{x\}$ to true. Disregarding the time to determine existence of violated clauses $C_{a'}$, this step takes time 
\begin{equation}\label{eq:rec1}
\sum_{w=0}^d \sum_{\substack{\text{weight-}w\\ \text{assignment }a'}} \sum_{x\in \vars(C_{a'})\setminus \ones(a')} T(n, k-(w+1)) \le \sum_{w=0}^d O(n^w) T(n, k-(w+1)). 
\end{equation}
To determine a violated clause $C_{a'}$ (if it exists) for all weight-($\le d$) assignments $a'$, we simply traverse each clause $C$, determine the at most $\sum_{w=0}^d {r \choose w}= O(1)$ weight-($\le d$) assignments violating $C$ and store $C$ as violated for each of these assignments (if no other clause is already stored). This step takes time $O(m) = O(n^r)$ in the beginning. 
%
%
After this preprocessing, it remains to consider $d$-robust assignments. To determine whether a $d$-robust assignment satisfies $\phi$, we define a formula $\phi_d$ that is satisfied only by satisfying assignments of $\phi$, and particularly by all $d$-robust satisfying assignments of $\phi$. To this end, let $F_d$ contain all assignments of weight at most $d$ that violate some clause $C$ of $\phi$, and define
\[ \phi_d := \bigwedge_{a \in F_d} \NAND(\ones(a)). \]
  \begin{lem}\label{lem:NANDrep}
  The constructed formula $\phi_d$ has the following properties:
\begin{enumerate}[label=(P\arabic*)]
  \item If $\confam$ is $\NAND_{d+1}$-avoiding, then any satisfying assignment $a$ of $\phi_d$ is a satisfying assignment of $\phi$.
  \item If $a$ is a $d$-robust satisfying assignment of $\phi$, then $a$ satisfies $\phi_d$.
    \end{enumerate}
  \end{lem}
  \begin{proof}
  To prove (P1), we will make use of the following property.
  \begin{prop}\label{prop:NANDd1}
  Let $\confam$ be a $\NAND_{d+1}$-avoiding family. Then if an assignment $a$ violates some clause $C$ (chosen from $\confam$), there is an assignment $a' \le a$ of weight at most $d$ that violates~$C$.
  \end{prop}
  \begin{proof}
  We prove the claim via induction on the weight $w$ of the clause $C$ under $a$. If $w \le d$, the claim trivially holds. To prove the inductive step, we may assume for contradiction that there is an assignment $a$ of weight $w\ge d+1$ violating some clause $C=f(\bar{x})$, but no assignment $a' \le a$ of weight at most $w-1$ violates $C$. We will show that $\NAND_{d+1}$ can be obtained as a restriction of $f$. 
  To this end, choose some set $S \subseteq \ones(a) \cap \vars(C)$ of size $d+1$ (which is possible as $w\ge d+1$), and partition $\vars(C)$ into $S$, $Z_1 \coloneqq (\ones(a) \cap \vars(C)) \setminus S$ and $Z_0 \coloneqq \vars(C) \setminus (S\cup Z_1)$. Observe that we have   
  \begin{align*}
  \begin{matrix}
f( \overbrace{y_1\dots y_{d+1}}^{S},& \overbrace{\entry{0}}^{Z_0}, & \overbrace{\entry{1}}^{Z_1}) & = & 1, & \quad \text{if $(y_1,\dots,y_{d+1})\ne (1,\dots, 1)$}\\
  \hline
f( \; \entry{1}\quad,& \entry{0}, & \entry{1}) & = & 0. & \quad\text{[since $a$ violates $C$]}
    \end{matrix}
    \end{align*}
    where the first line follows since no assignment $a' \le a$ of weight at most $w-1$ violates $C$, yielding a contradiction.
  
  \cqed \end{proof}
  To prove (P1), assume that an assignment $a$ violates some clause $C$ of $\phi$. Since $\confam$ is $\NAND_{d+1}$-avoiding, by Proposition~\ref{prop:NANDd1} there exists an assignment $a'\le a$ of weight at most $d$ such that $a'$ violates $C$. Thus, $\phi_d$ contains a clause $\NAND(\ones(a'))$, which is violated by $a$, as $a' \le a$.
%
  To prove (P2), assume for contradiction that a $d$-robust assignment $a$ satisfies $\phi$ but not $\phi_d$. Then there is some $a' \in F_d$ such that $\NAND(\ones(a'))$ is violated by $a$, i.e., $a' \le a$. As $a' \in F_d$, there must be a clause $C$ of $\phi$ that is violated by $a'\le a$, which proves that $a$ is not $d$-robust and thus yields a contradiction.
  \end{proof}
%
  Note that $\phi_d$ is a $\SAT(\confam')$ formula with constraint family $\confam' = \{ \NAND_j \mid 2 \le j \le d\}$ of arity~$d$. Thus, by Proposition~\ref{prop:aritybaseline}, we can determine satisfiability of $\phi_d$ in time $f(k) (n^{2d} + \THC{d}(n, k))\log n$. We obtain the following recurrence by combining~\eqref{eq:rec1}, the $O(m)$-time preprocessing to determine violated classes $C_{a'}$, and $f(k)(n^{2d} + \THC{d}(n,k))\log n$ to solve $\phi_d$:  
  \begin{equation}\label{eq:rec}
  T(n,k) =  \Oh(m) + f(k)(n^{2d} + \THC{d}(n,k))\log n + \sum_{w=0}^d O(n^w) T(n, k-(w+1))
  \end{equation}
%
  Assume that there are $\gamma \ge d/(d+1)$ and $c$ such that $\THC{d}(n,k) \le g(k)n^{\gamma k+c}$. We will show that  $T(n,k) = g'(k)\Oh(n^{\gamma k + c'})$ for any $c'>\max\{c, 2r\}$ and $g'(k) = f(k)g(k)$. 
%
  We prove the claim via induction on $k$. The base case is $k < c'$, in which case we can solve $\SAT(\confam)$ in time $f(k)(n^{2r} + \THC{r}(n,k))\log n = f(k)\Oh((n^{2r}+n^k)\log n) \le \Oh(n^{c'})$, satisfying the claim. Thus, let us assume that $k \ge c'$ and that the claim holds for all $k'\le k-1$. Using~\eqref{eq:rec}, we obtain
\begin{align*}
T(n,k) & \le \Oh(m) + f(k)(n^{2d} + g(k) n^{\gamma k + c})\log n + g'(k) \left(\sum_{w=0}^d O(n^w) n^{\gamma (k-w+1) + c'}\right) \\
    & \le g'(k) \log n\cdot \Oh\left(n^{2r} + n^{\gamma k + c} + \sum_{w=0}^d n^{w + \gamma(k-(w+1)) + c'}\right) \\
    & \le g'(k) \log n\cdot \Oh\left(n^{2r} + n^{\gamma k + c} + n^{\gamma k + c'}\right) = g'(k) \Oh(n^{\gamma k + c'}),
\end{align*}
where in the second line, we used that $g'(k) = f(k)g(k)$, and in the last line we used that $\gamma(w+1) \ge w$ as $\gamma \ge d/(d+1) \ge w/(w+1)$ for $w\le d$, as well as our choice of $c'$ which satisfies $c' > c$ and $\gamma k + c' \ge c' > 2r$.
%
%
\section{Hardness Results}
\label{sec:hardness}
%
In this section, we give our hardness results. To this end, we first consider $\Implication = \SAT(\Impl)$ and give a $f(k)n^{(\omega/6-o(1))\sqrt[3]{k}}$-lower bound under the $k$-clique conjecture. Afterwards, we handle the case of $\NAND_{d}$- or $\Impl$-representing families, by reducing from $d$-uniform (Hyper)Clique or \Implication, respectively.
%
\subsection{Hardness for \Implication}
%
\begin{thm}\label{thm:implhardness}
If \Implication can be solved in time $f(k) n^{(\omega/6-\eps)\sqrt[3]{k}+c}$ for some $\eps>0, c$ and $f(k)$, then the $k$-Clique conjecture fails.
\end{thm}
\begin{proof}
%
Let $G=(V, E)$ be an undirected graph.
We construct an \wDAGImplication instance $G' = (V', E', w')$ with parameter $k' = k\cdot K + {k \choose 2}$ with $K\coloneqq {k \choose 2}+1$ as follows.
The vertex set $V'$ is the disjoint union of \emph{vertex nodes} $V'_V \coloneqq \{v_u \mid u\in V\}$  and \emph{edge nodes} $V'_E \coloneqq \{v_e \mid e\in E\}$. For every $e=\{u,w\}\in E$, we introduce the edges $(v_e, v_u), (v_e, v_w)$ to $E'$. Furthermore, we set the weights of vertex nodes to $K$, and the weights of edge nodes to $1$. 
%
\begin{claim}
There is a closed set $X$ of weight $k'$ in $G'$ if and only if there is a $k$-clique in $G$.
\end{claim}
\begin{proof}
Let $C=\{v_1, \dots, v_k\}$ be a $k$-clique in $G$. Observe that $X = \{ v_u \mid u\in C\} \cup \{ v_e \mid e\in {C\choose 2}\}$ is a closed set in $G'$ of weight $|C|K + {|C| \choose 2} = k\cdot K + {k \choose 2} = k'$.
%
For the converse, assume that $X$ is a closed set in $G'$ of weight $k'$. Setting $X_V := X\cap V'_V$ and $X_E := X\cap V'_E$, we show the following sequence of facts:
\begin{enumerate}[label=\arabic*)]
\item $X_E \subseteq {X_V \choose 2}$: note that $X$ is only closed if for all $v_{\{u,w\}} \in X_E$, we have $v_u, v_w \in X_V$.
\item $|X_V| = k$ and $|X_E| = {k \choose 2}$: note that if $|X_V|< k$, then $|X_E| \le {k-1 \choose 2}$ by 1) and thus the weight of $X$ is $|X_V|K+|X_E| \le (k-1)K + {k-1 \choose 2} < kK+ {k \choose 2} = k'$. Furthermore, if $|X_V|>k$, then the weight of $X$ is at least $|X_V|K \ge (k+1)K = kK + {k \choose 2} + 1> k'$. Thus, we have $|X_V| = k$, and hence we must have $|X_E| = {k \choose 2}$ for $|X_V|K + |X_E| = k'$ to hold.
%
\item $X_V$ forms a $k$-clique in $G$: Facts 1) and 2) require that $X_E = {X_V \choose 2}$, which implies that $E$ contains all edges between vertices of $X_V$. 
\end{enumerate}
The last statement concludes the proof of the claim.
\cqed
\end{proof}
%
Assume that for some $c$ and $\eps >0$, there is an \Implication algorithm running in time $f(k)n^{(\omega/6-\eps)\sqrt[3]{k}+c}$. Given a $k$-clique instance $G$, we run the above reduction to create a \wDAGImplication instance $G'$ with parameter $k'\le (k+1)({k \choose 2} + 1) = (k^3 + k+2)/2 \le k^3$ for $k\ge 2$. Observe that $G'$ has $\Oh(n^2)$ nodes and can be converted to an equivalent \Implication instance $G''$ with the same parameter $k'$ and $\Oh(k^2n^2)$ nodes by simulating each node weight~$w$ by a cycle of $w$ nodes. Now, we determine whether  $G''$ has a closed set of weight $k' \le k^3$ using the \Implication algorithm and thus decide $k$-clique in time $f(k^3)\Oh((k^2n^2)^{(\omega/6-\eps) k + c})=f(k^3)k^{\Oh(k)}n^{(\omega/3-2\eps) k + 2c}$, refuting the $k$-Clique conjecture.
\end{proof}
%
%
\subsection{Hardness for $\SAT(\confam)$}
%
In this section, we give our hardness results for general constraint families $\confam$ by reducing from ($d$-uniform Hyper-)Clique either via the independent set problem or via \Implication.
%
To obtain these results, we frequently have to plug-in constant $0$s or $1$s to obtain our desired constraints. Technically, this is a non-trivial step, as we need to enforce some variables to be assigned fixed values without blowing up the number of variables or the weight of the desired solution. To facilitate our proofs, we first formalize the problem variant that allows us to plug-in constants freely.
\begin{defn}
Let $\confam$ be an arbitrary constraint family and $\Sigma\subseteq \{0,1\}$. The problem $\SAT_\Sigma(\confam)$ asks to determine whether a given formula $\phi$ with Boolean variables $x_1,\dots, x_n$ is has a satisfying assignment of weight $k$, where $\phi$ is a conjunction of $m$ constraints of the form $f(\mathbf{x})$, where $f:\{0,1\}^r \to \{0,1\}$ is a constraint function in $\confam$ and $\mathbf{x}$ is an $r$-tuple over $\{x_1,\dots, x_n\} \cup \Sigma$ (any variable or constant $c\in \Sigma$ may be used repeatedly). Note that $\SAT_\emptyset(\confam) = \SAT(\confam)$.
\end{defn}
%
Ideally, we would like to show that $\SAT_{\{0,1\}}(\confam)$ is equivalent to $\SAT(\confam)$. More specifically, we would like to employ reductions of the following form. 
\begin{defn}
Let $\confam$ be an arbitrary constraint family, and $\Sigma,\Sigma' \subseteq \{0,1\}$ be disjoint. We say that $\SAT_\Sigma(\confam)$ \emph{expresses} $\Sigma'$, if there is a constant $c$ such that the following holds: For any formula $\phi$ of $\SAT_{\Sigma\cup \Sigma'}(\confam)$ and parameter $k$, we can compute, in linear time, a formula $\phi'$ of $\SAT_\Sigma(\confam)$ with parameter $k':= k + c$ such that $\phi$ has a satisfying assignment of weight $k$ if and only if $\phi'$ has a satisfying assignment of weight $k'$.
\end{defn}
%
Indeed, for $0$-invalid $\confam$, we can show that $\SAT(\confam)$ expresses $\{0,1\}$ (this is straightforward and was already shown in~\cite{Marx05}). For $0$-valid $\confam$, however, expressing the constant $1$ in general appears impossible. To still give tight hardness results for $\confam$ whenever it represents a hard function $g$, we make use of a stronger notion that captures whether we can obtain $g$ already as a restriction that avoids the constant $1$.
Formally, let $f:\{0,1\}^r \to \{0,1\}, g:\{0,1\}^s\to \{0,1\}$ be arbitrary Boolean functions. We say that a function $f$ \emph{contains $g$ as a $0$-restriction} if $g$ is obtained from $f$ by replacing each argument of $f$ either by an argument of $g$ or the constant $0$, i.e., we can partition $[r]$ into $X_1,\dots, X_s, Z_0$ such that 
\[ g(x_1,\dots,x_s) = f(\overbrace{\entry{x_1}}^{X_1}, \dots, \overbrace{\entry{x_s}}^{X_s}, \overbrace{\entry{0}}^{Z_0}).\]
%
  
  Using careful constructions, we can prove the following central technical lemma.
%
\begin{lem}\label{lem:representingconstants}
Let $\confam$ be an arbitrary constraint family and let $g$ be $\Impl$ or $\NAND_d$ for some $d\ge 2$. If some $f\in \confam$ contains $g$ as a restriction, then $\SAT(\confam)$ expresses $\{0,1\}$, or $\SAT(\confam)$ expresses $0$ and $f$ contains $g$ already as a $0$-restriction.
\end{lem}
%
Postponing the proof of the above lemma to the Sections~\ref{sec:0invalid} and~\ref{sec:0valid}, we can give the proof of our hardness results.
%
\begin{thm}[Hardness for $\SAT(\confam)$]
Let $\confam$ be a constraint family.
\begin{enumerate}
\item If $\confam$ represents $\Impl$, then $\SAT(\confam)$ cannot be solved in time $f(k)\Oh(n^{(\omega/6-\eps)\sqrt[3]{k}+c})$ for any computable $f(k)$ and constants $c, \eps > 0$, unless the $k$-Clique conjecture fails.
\item If $\confam$ represents $\NAND_{2}$, then $\SAT(\confam)$ cannot be solved in time $f(k)\Oh(n^{(\omega/3-\eps)k+c})$ for any computable $f(k)$ and constants $c, \eps > 0$, unless the $k$-Clique conjecture fails. 
\item If $\confam$ represents $\NAND_{d}$ for some $d\ge 3$, then $\SAT(\confam)$ cannot be solved in time $f(k)\Oh(n^{(1-\eps)k+c})$ for any computable $f(k)$ and constants $c, \eps > 0$, unless the $d$-uniform HyperClique conjecture fails. 
\end{enumerate}
\end{thm}
\begin{proof}
First, we observe that \Implication reduces to $\SAT(\Impl)$ such that 
\begin{equation}\label{eq:impltoSAT}
T_\Implication(n,k)\le O(T_{\SAT(\Impl)}(n,k)).
\end{equation}
Indeed, given any directed graph $G=(V,E)$ with $V=\{v_1,\dots,v_n\}$, we define the formula $\phi$ with variables $x_1,\dots,x_n$ and the set of constraints obtained by including $x_i \Rightarrow x_j$ for all $(v_i,v_j)\in E$. Note that for any $S\subseteq [n]$, $\{v_i\}_{i\in S}$ is a valid set in $G$ iff $a_S$ is a satisfying assignment of $\phi$, yielding~\eqref{eq:impltoSAT}.
%
Similarly, we observe that the $d$-uniform HyperClique problem reduces to $\SAT(\NAND_d)$ such that 
\begin{equation}\label{eq:HCtoSAT}
\THC{d}(n,k)\le O(T_{\SAT(\NAND_d)}(n,k)).
\end{equation}
Indeed, given any $d$-uniform hypergraph $G=(V,E)$ with $|V|= n$, we define the formula $\phi$ with variables $x_1,\dots,x_n$ and the constraints obtained by including, for each distinct $v_{i_1},\dots,v_{i_d}\in V$ such that $(v_{i_1},\dots,v_{i_d}) \notin E$, the constraint $\NAND_d(x_{i_1}, \dots, x_{i_d})$. Observe that $(v_{i_1},\dots,v_{i_k})\in V^k$ is a hyperclique in $G$ iff the weight-$k$ assignment with $x_{i_\ell}=1$ for all $\ell \in [k]$ satisfies $\phi$, yielding~\eqref{eq:HCtoSAT}. 
%
It remains to show that whenever some $f\in \confam$ contains $g \in \{\Impl\} \cup \{\NAND_d \mid d\ge 2\}$ as a restriction, then there is a computable $f'(k)$ and constant $c'$ such that
\begin{equation}\label{eq:restToFam}
T_{\SAT(g)}(n,k) \le f'(k) \cdot T_{SAT(\confam)}(n,k+c').
\end{equation}
Indeed, if $\SAT(\confam)$ expresses $\{0,1\}$, then 
\[ T_{\SAT(g)}(n,k) \le \Oh(T_{\SAT_{\{0,1\}}(\confam)}(n,k)) \le f'(k) \Oh(T_{\SAT(\confam)}(n,k+c')). \]
Here the first inequality follows by replacing each occurrence of a constraint $g(x_{i_1},\dots,x_{i_d})$ of $\SAT(g)$ by the corresponding restriction $f(g_1(x_{i_1},\dots,x_{i_d}),\dots, g_r(x_{i_1},\dots,x_{i_d}))$ of $\SAT_{\{0,1\}}(\confam)$. The second inequality follows from the definition of $\SAT(\confam)$ expressing $\{0,1\}$.
%
In the other case, $\SAT(\confam)$ expresses only $0$, but $f$ contains $g$ already as a $0$-restriction. Then we have
\[ T_{\SAT(g)}(n,k) \le \Oh(T_{\SAT_{\{0\}}(\confam)}(n,k)) \le f'(k) \Oh(T_{\SAT(\confam)}(n,k+c')), \]
as replacing each occurrence of a constraint $g(x_{i_1},\dots,x_{i_d})$ of $\SAT(g)$ by the corresponding restriction $f(g_1(x_{i_1},\dots,x_{i_d}),\dots, g_r(x_{i_1},\dots,x_{i_d}))$ does not require the use of the constant $1$. The second inequality again follows from the definition of $\SAT(\confam)$ expressing $0$.
%
As a consequence, by \eqref{eq:impltoSAT} and \eqref{eq:restToFam}, a $f(k) \cdot O(n^{(\omega/6-\eps)\sqrt[3]{k}+c})$ $\SAT(\confam)$ algorithm for an $\Impl$-representing family $\confam$ would then give an \Implication algorithm running in time  
\[f(k) f'(k) \Oh(n^{(\omega/6-\eps)\sqrt[3]{k+c'}+c})= f''(k) \Oh(n^{(\omega/6-\eps)\sqrt[3]{k}+c''}),\]
where $f''(k) = f(k) f'(k)$ and $c'' \le c+\sqrt[3]{c'}$. This would refute the $k$-Clique conjecture by Theorem~\ref{thm:implhardness}, concluding 1.
%
Similarly, a $f(k) \cdot O(n^{\gamma k+c})$ $\SAT(\confam)$ algorithm for an $\NAND_d$-representing family $\confam$ would give a $d$-uniform HyperClique algorithm running in time  
\[f(k) f'(k) \Oh(n^{\gamma k+c+c'})= f''(k) \Oh(n^{\gamma k +c''}),\]
where $f''(k) = f(k) f'(k)$ and $c'' = c+c'$. This yields 2. and 3. by the $k$-Clique or $d$-uniform HyperClique conjecture, respectively.
\end{proof}
%
In the remainder of the section, we prove Lemma~\ref{lem:representingconstants}. We split the proof in two cases, depending on whether $f$ is 0-invalid (Lemma~\ref{lem:0invalidexpressesboth}) or 0-valid (Corollary~\ref{cor:0validfmainclaim}).
%
\subsection{Proof of Lemma~\ref{lem:representingconstants}: 0-invalid case}
\label{sec:0invalid}
%
%
Let $f$ be such that we can obtain $\Impl$ or $\NAND_d$ for $d\ge 2$ as a restriction. Note that if it contains $\NAND_{d},d>2$ then it also must contain $\NAND_2$ as a restriction. 
%
In this section, we consider the case that $f(y_1,\dots,y_r)$ is not $0$-valid, i.e., the all-zeroes assignment $u_1=\cdots = u_r= 0$ does not satisfy $f$.
%
\begin{lem}\label{lem:0invalidexpressesboth}
If $f$ contains $\Impl$ or $\NAND_2$ as a restriction and $f$ is 0-invalid, then $\SAT(\confam)$ expresses $\{0,1\}$.
\end{lem}
%
The above result in fact follows from the following claim.
%
\begin{claim}\label{claim:0invalidconst}
Let $f$ be as above. Given a parameter $k'$, we can compute, in time $\Oh(k')$, a formula $\phi_{0,1}$ of $\SAT(\confam)$ with variables $y, z_1, \dots, z_{k'+1}$ such that the only satisfying assignment of weight at most $k'$ is $y=1, z_1=\cdots = z_{k'+1} = 0$.
\end{claim}
%
Indeed, let us assume the above claim, and take any formula $\phi$ of $\SAT_{\{0,1\}}(\confam)$ with parameter~$k$. We construct $\phi_{0,1}$ with parameter $k':=k+1$ and define the formula $\phi'$ on variable set $x_1,\dots,x_n, y, z_1,\dots, z_{k'+1}$ where we include all constraints of $\phi_{0,1}$ and all constraints of $\phi$, \emph{replacing each use of the constant $0$ by $z_1$ and each use of the constant $1$ by $y$}. This yields a formula of $\SAT(\confam)$ with the property that for any weight-$k$ solution $x_1,\dots,x_n$ of $\phi$, the corresponding assignment that sets $y=1$ and $z_1=\cdots=z_{k'+1} = 0$ is a weight-$(k+1)$ solution of $\phi'$. Conversely, any $(k+1)$-weight solution of $\phi'$ must set $y=1$ and $z_1=0$ by the above claim, and hence the assignment to $x_1,\dots,x_n$ must also satisfy $\phi$. Observe that this proves Lemma~\ref{lem:0invalidexpressesboth}. 
%
\begin{proof}[Proof of Claim~\ref{claim:0invalidconst}]
%
We first give a set of constraints that enforces $y=1$. To this end, let $S\subseteq [r]$ be such that $a_S$ satisfies $f$; observe that $S$ exists and is non-empty (otherwise $f$ contains neither $\Impl$ nor $\NAND_{2}$ as a restriction). For each $j=1,\dots,k'+1$, define the constraint $C_j$ obtained by plugging in $y$ for each $u_i$ with $i\in S$ (i.e., all arguments set to $1$ under $a_S$), and $z_j$ for all other values. We claim that any weight-$(\le k')$ assignment satisfying $\bigwedge_{j=1}^{k'+1} C_j$ sets $y=1$: by the weight restriction, at least one of $z_{1},\dots,z_{k'+1}$ must be equal to $0$, say $z_{j^*}$. Then setting $y=0$ would falsify $C_{j^*}$, as then all its arguments are $0$. Note, however, that the desired assignment $y=1, z_1=\cdots=z_{k'+1}=0$ satisfies $\bigwedge_{j=1}^{k'+1} C_j$.
%
It remains to give additional constraints enforcing that $z_j = 0$ for \emph{all} $j\in [k'+1]$. As a first step, we find $S\subsetneq T$ such that $f(a_S)=1$ but $f(a_T)= 0$: Since $f$ represents $\Impl$ or $\NAND_2$, there is a partition of $[r]$ into $X, Y, Z_0, Z_1$ such that one of the following set of equalities hold:
\[
\begin{matrix}
f( \overbrace{\entry{0}}^{X},& \overbrace{\entry{0}}^{Y}, & \overbrace{\entry{0}}^{Z_0}, & \overbrace{\entry{1}}^{Z_1}) & = & 1,\\
f( \entry{0},& \entry{1}, & \entry{0}, & \entry{1}) & = & 1,\\
f( \entry{1},& \entry{1}, & \entry{0}, & \entry{1}) & = & 1,\\
  \hline
f( \entry{1},& \entry{0}, & \entry{0}, & \entry{1}) & = & 0.
\end{matrix}
\qquad 
\begin{matrix}
f( \overbrace{\entry{0}}^{X},& \overbrace{\entry{0}}^{Y}, & \overbrace{\entry{0}}^{Z_0}, & \overbrace{\entry{1}}^{Z_1}) & = & 1,\\
f( \entry{1},& \entry{0}, & \entry{0}, & \entry{1}) & = & 1,\\
f( \entry{0},& \entry{1}, & \entry{0}, & \entry{1}) & = & 1,\\
  \hline
f( \entry{1},& \entry{1}, & \entry{0}, & \entry{1}) & = & 0.
\end{matrix}
\]
In both cases, the first and fourth line yield sets $S \subsetneq T$ with $f(a_S)=1$ and $f(a_T)=0$ (specifically, for $S=Z_1$ and $T=X \cup Z_1$ or for $S=Z_1$ and $T=X \cup Y \cup Z_1$).
%
Given such $S,T$, for each $j,j' \in {[r] \choose 2}$, we define the constraint $C'_{j,j'}$ obtained from $f(u_1,\dots,u_r)$ by plugging-in $y$ for all $u_i$ with $i\in S$, $z_j$ for all $i\in T\setminus S$ and $z_{j'}$ for all other $i$. Note that any satisfying assignment of weight at most $k$ sets at least one of $z_1,\dots, z_{k'+1}$ to $0$, say $z_{j^*}$. Observe that the constraint $C'_{j,j^*}$ is satisfied iff $z_j=0$, as setting $z_j$ to 0 or 1 corresponds to the assignments $a_S$ (satisfying) or $a_T$ (unsatisfying), respectively. Furthermore, observe that setting $y=1$ and $z_1=\cdots = z_{k'+1}=0$ indeed satisfies all $C'_{j,j'}$.
This concludes the claim that the only satisfying assignment of weight at most $k'$ is $y=1, z_1=\cdots = z_{k'+1} = 0$.
\end{proof}
%
%
\subsection{Proof of Lemma~\ref{lem:representingconstants}: 0-valid case}
\label{sec:0valid}
%
In this section, we consider the case that $f(y_1,\dots,y_r)$ is $0$-valid, i.e., the all-zeroes assignment $u_1=\cdots = u_r= 0$ satisfies $f$. We first observe that we can still express at least the constant 0.
%
\begin{lem}\label{lem:0validexpresses0}
If some $f\in \confam$ contains $\Impl$ or $\NAND_2$ as a restriction and $f$ is 0-valid, then $\SAT(\confam)$ expresses $0$.
\end{lem}
\begin{proof}
Observe that it suffices to show how to construct, given a parameter $k$, a formula on variables $z_1,\dots, z_{k+1}$ such that the only satisfying assignment of weight at most $k$ sets $z_1=\cdots = z_{k+1}=0$.
%
To this end, assume first that $f$ is not satisfied by the all-ones assignment. Then, the formula $\bigwedge_{i=1}^{k+1} f(z_i,\dots,z_i)$ is trivially only satisfied by the assignment $z_1=\cdots = z_{k+1} = 0$.
%
Otherwise, observe that there must be a non-empty set $S\subsetneq [r]$ such that $a_S$ does not satisfy $f$ (otherwise $f$ would be a trivial constraint and could contain neither of $\Impl$ and $\NAND_2$). For each $i,i'\in [k+1]$, we define the constraint $C_{i,i'}$ obtained by using $z_i$ for all arguments in $S$, and $z_{i'}$ for all arguments not in $S$. Observe that $C_{i,i'} \wedge C_{i',i}$ forces $z_i = z_{i'}$, and thus $z_1= \cdots = z_{k+1}$, which is satisfied by an assignment of weight at most $k$ if and only if the common value is $0$. 
\end{proof}
%
Interestingly, for $0$-valid $f$, containing $\Impl$ as a restriction is equivalent to containing $\Impl$ already as a $0$-restriction.
%
\begin{lem}\label{lem:Implas0restr}
If $f$ contains $\Impl$ as a restriction and is $0$-valid, then $f$ contains $\Impl$ already as a $0$-restriction.
\end{lem}
\begin{proof}
Since $f:\{0,1\}^r \to \{0,1\}$ contains $\Impl$ as a restriction, we can partition $[r]$ into sets $X,Y,Z_0,Z_1$ and write
\begin{align}\label{eq:containsImpl}
\begin{matrix}
f( \overbrace{\entry{0}}^{X},& \overbrace{\entry{0}}^{Y}, & \overbrace{\entry{0}}^{Z_0}, & \overbrace{\entry{1}}^{Z_1}) & = & 1,\\
f( \entry{0},& \entry{1}, & \entry{0}, & \entry{1}) & = & 1,\\
f( \entry{1},& \entry{1}, & \entry{0}, & \entry{1}) & = & 1,\\
  \hline
f( \entry{1},& \entry{0}, & \entry{0}, & \entry{1}) & = & 0.\\
\end{matrix}
\end{align}
Assume first that 
\begin{equation}\label{eq:casedist}
f(\overbrace{\entry{0}}^{X}, \overbrace{\entry{1}}^Y, \overbrace{\entry{0}}^{Z_0}, \overbrace{\entry{0}}^{Z_1}) = 0.
\end{equation}
Then, we obtain $\Impl$ as a 0-restriction by setting $X'\coloneqq Y, Y'\coloneqq Z_1, Z' \coloneqq X\cup Z_0$ and observing that
\begin{align*}
\begin{matrix}
f( \overbrace{\entry{0}}^{X'=Y},& \overbrace{\entry{0}}^{Y'=Z_1}, & \overbrace{\entry{0}}^{Z'=X\cup Z_0}) & = & 1, & \quad \text{[$f$ is $0$-valid]}\\
f( \entry{0},& \entry{1}, & \entry{0}) & = & 1, & \quad\text{[by \eqref{eq:containsImpl}]}\\
f( \entry{1},& \entry{1}, & \entry{0}) & = & 1, & \quad\text{[by \eqref{eq:containsImpl}]}\\
  \hline
f( \entry{1},& \entry{0}, & \entry{0}) & = & 0. & \quad\text{[by \eqref{eq:casedist}]}\\
\end{matrix}
\end{align*}
Otherwise, if \eqref{eq:casedist} does not hold, then we obtain $\Impl$ as a 0-restriction by setting $X'\coloneqq X\cup Z_1, Y' \coloneqq Y, Z' \coloneqq Z_0$ and observing that 
\begin{align*}
\begin{matrix}
f( \overbrace{\entry{0}}^{X'=X\cup Z_1},& \overbrace{\entry{0}}^{Y'=Y}, & \overbrace{\entry{0}}^{Z'=Z_0}) & = & 1, & \quad \text{[$f$ is $0$-valid]}\\
f( \entry{0},& \entry{1}, & \entry{0}) & = & 1, & \quad\text{[by $\neg$\eqref{eq:casedist}]}\\
f( \entry{1},& \entry{1}, & \entry{0}) & = & 1, & \quad\text{[by \eqref{eq:containsImpl}]}\\
  \hline
f( \entry{1},& \entry{0}, & \entry{0}) & = & 0. & \quad\text{[by \eqref{eq:containsImpl}]}\\
\end{matrix}
\end{align*}
\end{proof}
%
%
It remains to handle the case that $f$ contains $\NAND_d$ as a restriction.
We first observe that if $f$ contains $\Impl$ as a $0$-restriction, then $\SAT_0(\confam)$ even expresses the constant $1$. (Thus, afterwards, we may assume that $f$ does not contain $\Impl$ as a 0-restriction.)
%
\begin{lem}\label{lem:0validexpresses1}
If some $f\in \confam$ contains $\Impl$ as a $0$-restriction, then $\SAT_0(\confam)$ expresses $1$. 
\end{lem}
\begin{proof}
Given any formula $\phi$ of $\SAT_{\{0,1\}}(\confam)$ on variables $x_1,\dots,x_n$, we construct a formula $\phi'$ on variables $x_1,\dots, x_n, y$ as follows: Since some $f\in \confam$ contains $\Impl$ as a $0$-restriction, we can express, for any variables $v,v'$, the implication $v \Rightarrow v'$ by a corresponding constraint of $\SAT_0(\confam)$. We construct $n$ such constraints to enforce $\bigwedge_{j=1}^{n} (x_j\Rightarrow  y)$. Subsequently, we may use $y$ to replace any use of the constant $1$ to convert the constraints of $\phi$ to constraints of the $\SAT_{0}(\confam)$-formula $\phi'$. 
%
To argue correctness, note that any satisfying weight-$k$ assignment of $\phi$ yields a satisfying weight-$(k+1)$ assignment of $\phi'$ by setting $y=1$. Conversely, note that any weight-$(k+1)$-assignment of $\phi'$ must set $y=1$ (since $k\ge 1$ implies that at least one variable $x_i$ is set to one, which enforces $y=1$ by the corresponding implication $x_i \Rightarrow y$) and thus corresponds to a weight-$k$ assignment to $x_1,\dots, x_n$ satisfying $\phi$. 
\end{proof}
%
%
%
%
%
In the remainder of this section, we assume that $f$ contains $\NAND_d$ as a restriction, but does not contain $\Impl$ as a $0$-restriction, and the aim is to find $\NAND_d$ already as a $0$-restriction.
%
%
\begin{lem}\label{lem:subsetdiff}
For any $0$-valid $f$, if $f$ does not contain $\Impl$ as a 0-restriction, then whenever $f(a_S) = f(a_T) = 1$ with $S\subseteq T$, then $f(a_{T\setminus S}) = 1$.
\end{lem}
\begin{proof}
If $S=T$, there is nothing to show, so let $S\subsetneq T$ and assume for contradiction that $f(a_{T\setminus S}) = 0$. We obtain $\Impl$ as a 0-restriction as follows:
\begin{align*}
\begin{matrix}
f( \overbrace{\entry{0}}^{X=T\setminus S},& \overbrace{\entry{0}}^{Y=S}, & \overbrace{\entry{0}}^{Z=[r]\setminus T}) & = & 1, & \quad \text{[$f$ is $0$-valid]}\\
f( \entry{0},& \entry{1}, & \entry{0}) & = & 1, & \quad\text{[$f(a_S)=1$]}\\
f( \entry{1},& \entry{1}, & \entry{0}) & = & 1, & \quad\text{[$f(a_T)=1$]}\\
  \hline
f( \entry{1},& \entry{0}, & \entry{0}) & = & 0. & \quad\text{[by assumption]}\\
\end{matrix}
\end{align*}
This yields the claim.
\end{proof}
%
We can finally obtain $\NAND_d$ as a $0$-restriction.
%
%
\begin{lem}\label{lem:NANDas0restr}
If $f$ contains $\NAND_d$ as a restriction, does not contain $\Impl$ as a 0-restriction and is $0$-valid, then $f$ contains $\NAND_d$ already as a $0$-restriction.
\end{lem}
\begin{proof}
%
%
%
%
%
%
%
%
%
%
%
%
%
%
%
%
%
%
%
%
%
%
%
%
%
%
%
%
%
%
%
%
%
%
%
%
%
%
%
%
%
%
%
%
Since $f:\{0,1\}^r \to \{0,1\}$ contains $\NAND_d$ as a restriction, we can partition $[r]$ into sets $X_1,\dots,X_d,Z_0,Z_1$ such that $X_I\coloneqq \bigcup_{i\in I} X_i$ with $I\subseteq [d]$ satisfies:
\begin{equation}\label{eq:NANDrestr}
f(a_{X_I \cup Z_1}) = \begin{cases} 0 & \text{if } I = [d], \\ 1 & \text{if } I \subsetneq [d]. \end{cases}
\end{equation}
%
We claim that the partition $X'_{i} \coloneqq X_i$ for $i<d$, $X'_d \coloneqq X_d \cup Z_1$, $Z'\coloneqq Z_0$ provides $\NAND_d$ as a 0-restriction: Letting $X'_I \coloneqq \bigcup_{i\in I} X'_i$, this follows from
\begin{equation}\label{eq:NAND0restr}
f(a_{X'_I}) = \begin{cases} 0 & \text{if } I = [d], \\ 1 & \text{if } I \subsetneq [d]. \end{cases}
\end{equation}
To verify~\eqref{eq:NAND0restr}, note first that $f(a_{X'_{[d]}}) = f(a_{X_{[d]}\cup Z_1}) = 0$ by~\eqref{eq:NANDrestr}. Second, let $I\subsetneq [d]$. If $d\in I$, then $f(a_{X'_I}) = f(a_{X_I\cup Z_1}) = 1$ by~\eqref{eq:NANDrestr}. Otherwise, if $d\notin I$, then we have $f(a_{X'_I}) = f(a_{X_I}) = 1$ by Lemma~\ref{lem:subsetdiff} (for this, note that $f$ does not contain $\Impl$ as 0-restriction and that $f(a_{X_I \cup Z_1}) = f(a_{Z_1}) = 1$). This concludes the claim.   
\end{proof}
%
The proof of this section is summarized in the following corollary. 
%
\begin{cor}\label{cor:0validfmainclaim}
If $f$ contains $g\in \{\Impl\} \cup \bigcup_{d\ge 2} \{\NAND_d\}$ and $f$ is $0$-valid, then $\SAT(f)$ expresses $\{0,1\}$, or $\SAT(f)$ expresses $0$ and contains $g$ as a $0$-restriction. 
\end{cor}
\begin{proof}
If $g=\Impl$, then $f$ contains $g$ already as a $0$-restriction by Lemma~\ref{lem:Implas0restr} and $\SAT(f)$ expresses $\{0,1\}$ by Lemmas~\ref{lem:0validexpresses0} and~\ref{lem:0validexpresses1}.
%
If $g=\NAND_d$, then either $f$ also contains $\Impl$ as a $0$-restriction, in which case $\SAT(f)$ expresses $\{0,1\}$ by Lemmas~\ref{lem:0validexpresses0} and~\ref{lem:0validexpresses1}, or it does not contain $\Impl$ as a $0$-restriction, and thus $f$ contains $g$ as a $0$-restriction by Lemma~\ref{lem:NANDas0restr} and $\SAT(f)$ expresses $0$ by Lemma~\ref{lem:0validexpresses0}.
\end{proof}
%
%
%
